%% kapa: exam=1 round=26 year=2015 period=1 form=A subject=civil_law subject_ko=민법 seq=3 number=3
%% source: Q-Net 감정평가사 기출문제 (cst00309) artlSeq=5138001 file=extracted/1008/5138001/제26회_감정평가사 1차_기출/제26회_감정평가사 1차_기출/1교시_감평_A형.pdf
%% conversion: kordoc 4.24.51 -> tools/gen1.py (auto); review=proofread-sonnet; flags=none
\begin{question}{3}
민법상 제한능력자에 관한 설명으로 옳지 않은 것은?
\begin{choices}
  \item 미성년자의 취소할 수 있는 단독행위는 추인이 있을 때까지 상대방이 거절할 수 있다.
  \item 피성년후견인이 일용품을 구입한 경우 성년후견인은 이를 취소할 수 없다.
  \item 가정법원은 성년후견개시의 심판뿐 아니라 한정후견개시의 심판을 할 때에도 본인의 의사를 고려하여야 한다.
  \item 특정후견은 본인의 의사에 반하여 할 수 없다.
  \item 가정법원이 피성년후견인에 대해 한정후견개시의 심판을 할 때에는 종전의 성년후견의 종료 심판을 할 필요가 없다.
\end{choices}
\end{question}
\answer{5}
